%!TeX program = xelatex

\documentclass[aspectratio=169]{beamer}
\usepackage[english]{babel}
\usepackage{fontspec}
\usepackage{graphicx}
\usepackage{xcolor}
\usepackage{xltxtra}
\usepackage{hyperref}
\usepackage{multicol}
\usepackage{unicode-math}
\usepackage{changepage}
\usepackage{tcolorbox}

\usetheme{Warsaw}
\setbeamertemplate{navigation symbols}{}
\setbeamertemplate{headline}{}
\setbeamertemplate{bibliography item}[text]

\definecolor{MainRed}{RGB}{175,44,34}     % CMYK 3-95-94-29
\definecolor{DarkRed}{RGB}{127,37,2}      % CMYK 22-91-82-46
\definecolor{AltRed}{RGB}{150,40,14}

\definecolor{GoldLight}{RGB}{200,160,110} % CMYK 18-35-60-10
\definecolor{GoldDark}{RGB}{150,121,82}   % CMYK 18-35-60-40

\definecolor{GrayOne}{RGB}{198,198,198}
\definecolor{GrayTwo}{RGB}{160,160,160}
\definecolor{GrayThr}{RGB}{135,135,135}

\definecolor{NearBlack}{RGB}{31,31,31}

\definecolor{Bejevi}{RGB}{201,188,171}
\definecolor{NearWhite}{RGB}{227,227,227}

\setbeamercolor{structure}{fg=MainRed}
\setbeamercolor{palette primary}{fg=white,bg=DarkRed}
\setbeamercolor{palette secondary}{fg=white,bg=MainRed}
\setbeamercolor{palette tertiary}{fg=white,bg=DarkRed}
\setbeamercolor{palette quaternary}{fg=white,bg=MainRed}

\setbeamercolor{title}{fg=white,bg=MainRed}
\setbeamercolor{frametitle}{fg=white,bg=MainRed}
\setbeamercolor{titlelike}{fg=white,bg=MainRed}

\setbeamercolor{normal text}{fg=black,bg=white}

\setbeamercolor{block title}{fg=white,bg=DarkRed}
\setbeamercolor{block body}{fg=black,bg=DarkRed!8}

\setbeamercolor{block title alerted}{fg=white,bg=MainRed}
\setbeamercolor{block body alerted}{fg=black,bg=MainRed!10}

\setbeamercolor{block title example}{fg=white,bg=GoldDark}
\setbeamercolor{block body example}{fg=black,bg=GoldLight!20}

\setbeamercolor{title in head/foot}{fg=white,bg=DarkRed}
\setbeamercolor{author in head/foot}{fg=white,bg=NearBlack}
\setbeamercolor{date in head/foot}{fg=white,bg=GrayTwo}

\setbeamertemplate{footline}{
        \leavevmode%
        \hbox{%
                \begin{beamercolorbox}[wd=.72\paperwidth,ht=2.8ex,dp=1.2ex,leftskip=1em]{title
                        in head/foot}%
                        \insertshorttitle
                \end{beamercolorbox}%
                \begin{beamercolorbox}[wd=.14\paperwidth,ht=2.8ex,dp=1.2ex,center]{author
                        in head/foot}%
                        Шаго П. Е.
                \end{beamercolorbox}%
                \begin{beamercolorbox}[wd=.14\paperwidth,ht=2.8ex,dp=1.2ex,center]{date
                        in head/foot}%
                        \insertframenumber/\inserttotalframenumber
                \end{beamercolorbox}%
        }%
        \vskip0pt%
}

\tcbuselibrary{skins}
\newtcolorbox{shadowbox}{
        enhanced,
        colback=GrayOne,
        colframe=white,
        boxrule=0pt,
        sharp corners=all,
        arc=6pt,
        drop shadow,
        left=4mm,
        right=4mm,
        top=2mm,
        bottom=2mm
}

\newtcolorbox{shadowbox2}{
        enhanced,
        colback=NearWhite,
        colframe=white,
        boxrule=0pt,
        arc=4pt,
        outer arc=4pt,
        fuzzy shadow={
        0mm}{-1.3mm}{0mm}{0.3mm}{black!30
        },
        coltext=NearBlack,
        halign=center,
        fontupper=\centering,
        left=4mm,
        right=4mm,
        top=2mm,
        bottom=2mm
}

\newtcolorbox{shadowbox3}{
        enhanced,
        colback=NearWhite,
        colframe=MainRed,
        coltext=NearBlack,
        arc=6pt,
        boxrule=0.6pt,
        left=4mm, right=4mm, top=2mm, bottom=2mm,
        halign=center
}

\setmainfont{Iosevka NFM}
\setmonofont{Iosevka NFM}
\setsansfont{Iosevka NFM}
\setromanfont{Iosevka NFM Italic}
\setmathfont{Latin Modern Math}

\title{Математическое моделирование и
        реализация планировщика
процессов для гетерогенных процессорных архитектур}

\begin{document}

\begin{frame}[plain]
        \begin{center}
                \begin{shadowbox3}
                        {\scriptsize
                                Санкт-Петербургский государственный
                                университет \\ [0.1em]
                                Факультет прикладной математики~-- процессов
                                управления \\ [-0.4em]
                                ООП СВ.5005.2022 Прикладная
                                математика, фундаментальная информатика и
                                программирование
                        }
                \end{shadowbox3}
        \end{center}

        \vspace{0.3em}
        {\color{MainRed}\hrule height 1pt}
        \vspace{0.3em}

        \begin{center}
                {\small\color{MainRed}\textbf{Шаго Павел Евгеньевич}}

                \vspace{0.2em}
                \begin{shadowbox3}
                        {\bfseries
                                Математическое моделирование и
                                реализация планировщика
                        процессов для гетерогенных процессорных архитектур}
                \end{shadowbox3}

                \vspace{0.1em}
                {\small\itshape Бакалаврская работа}
        \end{center}
        \vspace{0.5em}

        \begin{minipage}{\textwidth}
                \begin{minipage}[t]{0.45\linewidth}
                        \raggedright
                        {\scriptsize
                                \textit{Научный руководитель:}\\[0.15em]
                                кандидат физ.-мат. наук, доцент\\
                        Корхов В.\,В.}
                \end{minipage}
                \hfill
                \begin{minipage}[t]{0.45\linewidth}
                        \raggedright
                        {\scriptsize
                                \textit{Рецензент:}\\[0.15em]
                                эксперт по архитектуре и сетевому стеку,\\
                        ООО НИЦ АЭРОДИСК, Анохин Ю.\,В.}
                \end{minipage}
        \end{minipage}

        \vfill
        {\color{MainRed}\hrule height 1pt}

        \begin{center}
                {\scriptsize
                        Санкт-Петербург\\[-0.5em]
                2026}
        \end{center}
\end{frame}

\begin{frame}{Постановка задачи}
        \begin{columns}[T]
                \column{0.5\textwidth}
                \begin{itemize}
                        \item Современные процессоры всё чаще имеют гетерогенную
                                компоновку, в которой
                                производительные P\nobreakdash-ядра дополняются
                                энергоэффективными E\nobreakdash-ядрами

                        \item Стандартный планировщик Linux не
                                различает типы ядер,
                                поэтому не всегда способен
                                сбалансировать производительность
                                и энергопотребление так, как это
                                задумано архитектурой CPU
                \end{itemize}

                \column{0.5\textwidth}
                \textbf{Цели работы}
                \begin{enumerate}
                \item Формализовать распределение временных слотов
                        P/E\nobreakdash-ядер
                        как динамический аукцион Викри (VCG), в котором
                        задачи конкурируют за вычислительный ресурс

                \item Реализовать полученный механизм используя
                        \texttt{sched\_ext} и оценить
                        его по метрикам на тестах производительности
                \end{enumerate}
        \end{columns}

        \vspace{0.2em}
        \begin{shadowbox2}
                \small Каждое ядро рассматривается как
                \textbf{истощаемое благо}, неиспользованный квант
                времени теряется безвозвратно. Задача\nobreakdash-агент
                характеризуется типом $\theta_i = (v_i, l_i)$, а
                механизм максимизирует социальное благосостояние.
        \end{shadowbox2}
\end{frame}

\begin{frame}{Модель системы}
        \begin{columns}[T]
                \column{0.5\textwidth}
                \textbf{Аппаратная платформа}
                \begin{itemize}
                        \item Ядра: $\mathcal{P} = \{P_1,\ldots,P_{K_P}\}$,\;
                                $\mathcal{E} = \{E_1,\ldots,E_{K_E}\}$
                        \item Дискретное время $t = 1,2,3,\ldots$
                        \item Замедление на E-ядре:
                                $\sigma = \eta_P/\eta_E > 1$
                        \item Общий ISA (любая задача может быть
                                выполнена на любом ядре)
                \end{itemize}

                \column{0.5\textwidth}
                \textbf{Задачи как агенты}
                \begin{itemize}
                        \item Тип $\theta_i = (v_i, l_i)$ (оценка и длина)
                        \item Полезность $u_i = v_i - p_i$ при полном
                                выполнении
                        \item Бюджет $B_i$ по классу обслуживания:
                \end{itemize}
                \vspace{0.2em}
                {\footnotesize
                \begin{tabular}{@{}ll@{}}
                        Реального времени & $B_{RT} \gg 1$ \\
                        Интерактивный     & $B_{int}$ \\
                        Пакетный          & $B_{batch}$ \\
                        Фоновый           & $B_{bg} \approx 0$
                \end{tabular}}
        \end{columns}

        \vspace{0.4em}
        \begin{shadowbox}
                \footnotesize На E-ядре задача длиной $l_i$ исполняется за
                $\lceil l_i \cdot \sigma \rceil$ квантов.
                Ценность $v_i$ начисляется только при полном выполнении
                (аналог контракта Сано).
        \end{shadowbox}
\end{frame}

\begin{frame}{MDP модель}
        \begin{columns}[T]
                \column{0.5\textwidth}
                \textbf{Состояние} $s^t =
                (\mathbf{x}^t,\,\mathbf{b}^t,\,\boldsymbol{\omega}^t)$
                \begin{itemize}
                        \item $\mathbf{x}^t$~-- оставшиеся кванты на
                                каждом ядре
                        \item $\mathbf{b}^t$~-- остатки бюджетов активных задач
                        \item $\boldsymbol{\omega}^t$~-- внешние факторы
                                (температура, нагрузка)
                \end{itemize}

                \vspace{0.3em}
                \textbf{Действие} $a^t = (a_{i,P}^t, a_{i,E}^t)_i$
                при ограничениях:
                $$a_{i,P}^t + a_{i,E}^t \leq 1, \quad \forall i$$
                $$\sum_i a_{i,P}^t \leq K_P^t, \quad
                \sum_i a_{i,E}^t \leq K_E^t$$

                \column{0.5\textwidth}
                \textbf{Вознаграждение}

                \vspace{0.2em}
                Задача $i$ (покупатель):
                $$r_i(s,a) = v_i \cdot \mathbb{1}[\text{$i$ получила ядро}]$$

                Планировщик (продавец):
                $$r_0 = \mu_P\!\!\sum_{k\in\mathcal{P}}\!\mathbb{1}[x_k \geq 1]
                + \mu_E\!\!\sum_{k\in\mathcal{E}}\!\mathbb{1}[x_k \geq 1]$$

                \vspace{0.2em}
                \textbf{Социальное благосостояние}
                $$W(\pi) = \lim_{T\to\infty} \frac{1}{T}
                \mathbb{E}\!\left[\sum_{t=1}^T R(s^t,a^t)\right]$$
        \end{columns}
\end{frame}

\begin{frame}{Механизм VCG}
        \textbf{Эффективная ценность} задачи $i$ на ядре типа
        $\kappa \in \{P, E\}$
        $$\phi_P(\theta_i) = v_i - c_P \cdot l_i, \quad
        \phi_E(\theta_i) = v_i - c_E \cdot \lceil l_i \sigma \rceil$$

        \vspace{0.1em}
        \textbf{Оптимальное распределение} (максимизация $W$)
        $$a^t = \arg\max_{a} \sum_i \bigl[a_{i,P}\,\phi_P(\theta_i)
        + a_{i,E}\,\phi_E(\theta_i)\bigr]
        + \delta\, \mathbb{E}[W(s^{t+1})]$$

        \vspace{0.1em}
        \begin{columns}[T]
                \column{0.52\textwidth}
                \textbf{Платёж VCG} задачи $i$
                $$p_i^t = W_{-i}(s^t) - \bigl(W(s^t) - \phi_\kappa(\theta_i)\bigr)$$

                \vspace{0.1em}
                {\footnotesize Задача $i$ платит ровно то, что теряют
                остальные агенты из-за её присутствия (pivot rule)}

                \column{0.48\textwidth}
                \begin{shadowbox}
                        \footnotesize
                        \textbf{Свойства механизма}
                        \centering
                        \begin{itemize}
                                \item Стимульная совместимость
                                        (выгодно сообщать истинный $\theta_i$)
                                \item Индивидуальная рациональность
                                        $u_i \geq 0$
                                \item Бюджетное ограничение
                                        $p_i \leq B_i^t$
                        \end{itemize}
                \end{shadowbox}
        \end{columns}
\end{frame}

\begin{frame}{Маршрутизация задач}
        \textbf{Прокси для распределения}: отношение бюджета
        $\rho_i = B_i / B_i^{\max}$

        \vspace{0.3em}
        \begin{columns}[T]
                \column{0.33\textwidth}
                \begin{shadowbox2}
                        \small\textbf{DSQ\_P}

                        \vspace{0.3em}
                        \footnotesize
                        $\rho_i \geq \rho^*\ (50\%)$

                        \vspace{0.2em}
                        Интерактивные,\\«bursty» задачи

                        \vspace{0.2em}
                        Высокий $\rho_i$ = много\\простоя =
                        чувствитель-\\ность к задержке
                \end{shadowbox2}

                \column{0.33\textwidth}
                \begin{shadowbox2}
                        \small\textbf{DSQ\_E}

                        \vspace{0.3em}
                        \footnotesize
                        $\rho_i < \rho^*$

                        \vspace{0.2em}
                        CPU-bound,\\пакетные задачи

                        \vspace{0.2em}
                        Низкий $\rho_i$ = активно\\работает = подходит
                        \\для E-ядра
                \end{shadowbox2}

                \column{0.33\textwidth}
                \begin{shadowbox2}
                        \small\textbf{DSQ\_Starved}

                        \vspace{0.3em}
                        \footnotesize
                        $B_i < B_i^{\max} / 10$

                        \vspace{0.2em}
                        Бюджет исчерпан

                        \vspace{0.2em}
                        $v_d \mathrel{+}=
                        \max(0,\,-\phi_\kappa)\,/\,D$
                \end{shadowbox2}
        \end{columns}

        \vspace{0.3em}
        \begin{shadowbox}
                \footnotesize Бюджет пополняется пропорционально времени
                простоя: $B_i \mathrel{+}= \Delta t_{\text{idle}} \cdot v_i
                / D_{\text{rep}}$~--- аналог порогового правила в
                одномерных аукционах
        \end{shadowbox}
\end{frame}

\begin{frame}{Реализация (scx\_auction)}
        \begin{columns}[T]
                \column{0.57\textwidth}
                \textbf{Отклонения от абстрактной модели:}

                \vspace{0.2em}
                \textbf{(1) Непрерывная формула $\phi$} (нс, не кванты):
                $$\phi_P = v_i \cdot s - c_P \cdot l_{ns}, \qquad
                \phi_E = v_i \cdot s - c_E \cdot l_{ns} \cdot \sigma$$
                где $l_{ns}$~--- EWMA измеренного времени исполнения;
                $\phi < 0$ означает прямое направление в DSQ\_Starved

                \vspace{0.3em}
                \textbf{(2) Пропорциональный платёж}:
                $$p_i = c_\kappa \cdot t_{\text{consumed}}\,/\,s$$
                устраняет штраф за первый квант у коротких задач

                \column{0.43\textwidth}
                \begin{shadowbox}
                        \small\textbf{Параметры}

                        \vspace{0.2em}
                        \footnotesize
                        \begin{tabular}{ll}
                                \texttt{C\_P\_DEF}      & 512 \\[2pt]
                                \texttt{C\_E\_DEF}      & 256 \\[2pt]
                                \texttt{BURST\_PCT}     & 50\% \\[2pt]
                                \texttt{BUDGET\_MUL}    & 2000 \\[2pt]
                                \texttt{STARVE\_FRAC}   & 10 \\[2pt]
                                \texttt{P\_CAP\_PCT}    & 90\%
                        \end{tabular}
                \end{shadowbox}
        \end{columns}
\end{frame}

\begin{frame}{Среда тестирования}
        \begin{columns}[T]
                \column{0.24\textwidth}
                \begin{shadowbox}
                        \small\textbf{Система}

                        \vspace{0.25cm}
                        \footnotesize
                        Ядро:\\Linux 7.0

                        \vspace{0.25cm}
                        CPU:\\Intel Core\\ Ultra 7 265K\\
                        12E + 8P
                \end{shadowbox}

                \vspace{0.1cm}
                \begin{shadowbox}
                        \small\textbf{Сравнение}

                        \vspace{0.2cm}
                        \footnotesize
                        Linux EEVDF\\
                        \textit{vs}\\
                        scx\_auction (VCG)
                \end{shadowbox}

                \column{0.52\textwidth}
                \begin{shadowbox}
                        \small\textbf{Бенчмарки}

                        \vspace{0.1cm}
                        \footnotesize
                \begin{itemize}\setlength{\itemsep}{0.05cm}\setlength{\parskip}{0pt}
                                \item \textbf{hackbench}~--- 1000 задач\\
                                        Интенсивный обмен сообщениями,
                                        высокая нагрузка на планировщик.
                                        Оценивает переключение контекста.

                                \item \textbf{sysbench}~---
                                        стандартные параметры\\
                                        CPU, память, дисковый I/O.
                                        Моделирует OLTP под
                                        реалистичной нагрузкой.

                                \item \textbf{sched-latency}~---
                                        собственный eBPF\\
                                        Трейсинг трейспоинтов, запись
                                        p99 задержки планирования.
                        \end{itemize}
                \end{shadowbox}

                \column{0.24\textwidth}
                \begin{shadowbox}
                        \small\textbf{Метрики}

                        \vspace{0.2cm}
                        \footnotesize
                        \textbf{hackbench}\\
                        Время ($s$)

                        \vspace{0.2cm}
                        \textbf{sysbench}\\
                        $\text{Events}/s$

                        \vspace{0.2cm}
                        \textbf{sched-latency}\\
                        p99 ($\mu s$)
                \end{shadowbox}

        \end{columns}
\end{frame}

\begin{frame}{Результаты: Hackbench}
        \begin{center}
                \vspace{0.5em}
                \begin{shadowbox3}
                        {\Large\textbf{Нет данных}}\\[0.4em]
                        {\footnotesize Замеры scx\_auction в процессе сбора}
                \end{shadowbox3}

                \vspace{1em}
                \begin{shadowbox}
                        \footnotesize
                        \textbf{Метрика:} полное время выполнения ($s$),
                        меньше~--- лучше\\[0.3em]
                        \textbf{Нагрузка:} 1000 задач, интенсивный обмен
                        сообщениями\\[0.3em]
                        \textbf{Ожидание:} нейтральный результат~---
                        VCG-маршрутизация при высокой конкуренции
                        сводится к базовому EEVDF
                \end{shadowbox}
        \end{center}
\end{frame}

\begin{frame}{Результаты: Sysbench}
        \begin{center}
                \vspace{0.5em}
                \begin{shadowbox3}
                        {\Large\textbf{Нет данных}}\\[0.4em]
                        {\footnotesize Замеры scx\_auction в процессе сбора}
                \end{shadowbox3}

                \vspace{1em}
                \begin{shadowbox}
                        \footnotesize
                        \textbf{Метрика:} пропускная способность
                        ($\text{events}/s$), больше~--- лучше\\[0.3em]
                        \textbf{Нагрузка:} CPU-bound, однопоточная
                        вычислительная задача\\[0.3em]
                        \textbf{Ожидание:} задача с низким $\rho_i$
                        маршрутизируется на E-ядро; потенциальный
                        выигрыш при длинном фиксированном кванте
                \end{shadowbox}
        \end{center}
\end{frame}

\begin{frame}{Результаты: Задержка планирования}
        \begin{center}
                \vspace{0.5em}
                \begin{shadowbox3}
                        {\Large\textbf{Нет данных}}\\[0.4em]
                        {\footnotesize Замеры scx\_auction в процессе сбора}
                \end{shadowbox3}

                \vspace{1em}
                \begin{shadowbox}
                        \footnotesize
                        \textbf{Метрика:} p99 задержка планирования
                        ($\mu s$), меньше~--- лучше\\[0.3em]
                        \textbf{Нагрузка:} трейсинг трейспоинтов
                        под смешанной нагрузкой\\[0.3em]
                        \textbf{Ожидание:} интерактивные задачи
                        (высокий $\rho_i$) на P-ядре~---
                        потенциальное снижение хвостовой задержки
                        относительно однородного планировщика
                \end{shadowbox}
        \end{center}
\end{frame}

\begin{frame}{Анализ}
        \begin{center}
                \vspace{0.5em}
                \begin{shadowbox3}
                        {\Large\textbf{Нет данных}}\\[0.4em]
                        {\footnotesize Анализ будет проведён после
                        сбора результатов}
                \end{shadowbox3}

                \vspace{1em}
                \begin{shadowbox}
                        \footnotesize
                        \textbf{Ключевые вопросы:}
                        \begin{itemize}
                                \item Корректно ли $\rho_i$ разделяет
                                        интерактивные и CPU-bound задачи?
                                \item Насколько велик overhead VCG-платежей
                                        относительно ядерного EEVDF?
                                \item Подтверждается ли снижение хвостовой
                                        задержки для P-ядерных задач?
                        \end{itemize}
                \end{shadowbox}
        \end{center}
\end{frame}

\begin{frame}{Заключение}
        \begin{itemize}
                \item Задача распределения P/E-ядер формализована
                        как динамический VCG-аукцион в рамках MDP
                        \vspace{0.2cm}
                \item Механизм реализован через sched\_ext (eBPF)~---
                        политика изменяется без пересборки ядра
                        \vspace{0.2cm}
                \item Маршрутизация через отношение бюджета $\rho_i$~---
                        практическое приближение к оптимальному
                        пороговому правилу аукциона
                        \vspace{0.2cm}
                \item Направления развития: онлайн-обучение (IHMDP-VCG),
                        сравнение с ядерным EEVDF и s3+ (capacity-aware)
        \end{itemize}
\end{frame}

\begin{frame}{Список литературы \& Вопросы}
        \begin{thebibliography}{20}

                \bibitem{leon}
                Leon~G., Etesami~O.
                Mechanism Design for Scheduling with Heterogeneous Workers.
                Advances in Neural Information Processing Systems, 2022

                \bibitem{sano}
                Sano~T.
                Dynamic Mechanism Design for Multi-Period Contracting.
                Working paper, 2023

                \bibitem{sched_ext}
                \href{https://lwn.net/Articles/922405}
                {The extensible scheduler class (LWN)}

                \bibitem{eevdf}
                \href{https://lwn.net/Articles/925371}
                {An EEVDF CPU scheduler for Linux (LWN)}

        \end{thebibliography}
\end{frame}

\begin{frame}{Extra}
        \textbf{Стимульная совместимость (Sano, Теорема~1)}

        \vspace{0.2cm}
        Механизм стимульно совместим тогда и только тогда, когда
        для каждой задачи $i$:

        \begin{enumerate}
                \item $\alpha_i(v_i, l_i)$ не убывает по $v_i$
                        при фиксированном $l_i$
                \item $\int_0^{v_i} \alpha_i(\nu, l_i)\,d\nu$ не возрастает
                        по $l_i$ при фиксированном $v_i$
                \item Существует $\underline{\Pi}_i$, не зависящее от $l_i$:
                        $$\Pi_i(v_i, l_i) = \underline{\Pi}_i +
                        \int_0^{v_i} \alpha_i(\nu, l_i)\,d\nu$$
        \end{enumerate}

        \vspace{0.1cm}
        \begin{shadowbox}
                \footnotesize Задачам невыгодно завышать оценку $v_i$
                или занижать длину $l_i$~--- правдивая ставка
                является доминирующей стратегией
        \end{shadowbox}
\end{frame}
\addtocounter{framenumber}{-1}

\begin{frame}{Extra}
        \textbf{Онлайн-обучение: IHMDP-VCG (Leon \& Etesami)}

        \vspace{0.15cm}
        При неизвестных распределениях типов задач и ядре MDP планировщик
        обучается онлайн. Эпизодическая структура:

        \begin{itemize}
                \item \textbf{Фаза перемешивания} (длина $d^{[k]}$):
                        вывод цепи Маркова на стационарное распределение
                \item \textbf{Стационарная фаза} (длина $l^{[k]}$):
                        сбор платежей $\hat{p}^{[k]}$, обновление оценок
        \end{itemize}

        \vspace{0.15cm}
        \textbf{Гарантии} при $T \to \infty$ (с вер. $\geq 1 - O(\zeta)$):

        \vspace{0.1cm}
        \begin{columns}
                \column{0.5\textwidth}
                \begin{tabular}{ll}
                        Регрет благосостояния &
                        $\tilde{O}(\varepsilon T n + T^{2/3} n)$ \\[3pt]
                        Регрет продавца &
                        $\tilde{O}(\varepsilon T n^2 + T^{2/3} n^2)$ \\[3pt]
                        Регрет покупателей &
                        $\tilde{O}(\varepsilon T n^2 + T^{2/3} n^2)$
                \end{tabular}

                \column{0.5\textwidth}
                \begin{shadowbox}
                        \footnotesize Асимптотически: эффективность,
                        стимульная совместимость и индивидуальная
                        рациональность~--- одновременно
                \end{shadowbox}
        \end{columns}
\end{frame}
\addtocounter{framenumber}{-1}

\begin{frame}{Extra}
        \textbf{Пример: $K_P = K_E = 1$, одно свободное ядро типа $\kappa$}

        \vspace{0.2cm}
        \textbf{Победитель:}
        $$i^* = \arg\max_{i \in \mathcal{N}^t,\; B_i^t > 0}\;
        \phi_\kappa(\theta_i) + \delta^{m_\kappa(l_i)}\,\bar{W}_\kappa$$

        где $m_P(l) = l$,\; $m_E(l) = \lceil l\cdot\sigma \rceil$

        \vspace{0.2cm}
        \textbf{Платёж победителя:}
        $$p_{i^*} = \phi_\kappa^{-1}\!\left(\phi_\kappa(\theta_j) +
                \bigl(\delta^{m_\kappa(l_j)} - \delta^{m_\kappa(l_{i^*})}\bigr)
        \bar{W}_\kappa \;\Big|\; l_{i^*}\right)$$

        где $j$~--- задача на втором месте.

        \vspace{0.2cm}
        \begin{shadowbox}
                \footnotesize
                $\phi_P(\theta_i) > \phi_E(\theta_i)$ выполняется не всегда:
                при малом $v_i$ и большом $l_i$ E-ядро может дать
                большую эффективную ценность ($c_E \cdot \lceil l_i\sigma\rceil
                < c_P \cdot l_i$ при $\sigma < \gamma$).
                Оптимальная стратегия нетривиальна.
        \end{shadowbox}
\end{frame}
\addtocounter{framenumber}{-1}

\end{document}
